% Answer keys for Chapters 6--15. This file is \input directly inside the document.
% Every number comes from the check scripts in the temporary working directory.

\section*{Chapter 6: Vector Spaces and Subspaces}
\begin{enumerate}[label=\textbf{6.\arabic*}]
  \item $W$ is the solution set of the homogeneous system $x-2y+z=0$, $x+y=0$, so it is
        closed under addition and scalar multiplication. From $x=-y$ and $z=3y$ we get
        $W=\operatorname{span}\{(-1,1,3)\}$: a one-vector basis, $\dim W=1$.

  \item The set $D=\{A:a_{11}=a_{22}=0\}$ is closed under addition and scalars,
        because both diagonal entries stay zero. Every one of its members is
        $\begin{pmatrix}0&b\\c&0\end{pmatrix}=bE_{12}+cE_{21}$, so
        $\{E_{12},E_{21}\}$ is a basis and $\dim D=2$.

  \item $(1+x)+(1+x^{2})=2+x+x^{2}$: the set is linearly dependent, hence
        $\dim\operatorname{span}S=2$. The vector $x$ is not in that span, because
        $a(1+x)+b(1+x^{2})$ demands $a=1$ and $a+b=0$ at the same time. So
        $\{1+x,\ 1+x^{2},\ x\}$ is a basis of $P_{2}$.

  \item Collect the coefficients: $(c_{1}+c_{3})\mathbf{v}_{1}+(c_{1}+c_{2})
        \mathbf{v}_{2}+(c_{2}+c_{3})\mathbf{v}_{3}=\mathbf{0}$. Independence of
        $\{\mathbf{v}_{1},\mathbf{v}_{2},\mathbf{v}_{3}\}$ gives three homogeneous
        equations with determinant $2\ne0$, so $c_{1}=c_{2}=c_{3}=0$.

  \item $U$ is the $xy$-plane, $W$ is the plane $x=y$. Their intersection is
        $U\cap W=\operatorname{span}\{(1,1,0)\}$, so $\dim(U\cap W)=1$ and
        $\dim(U+W)=2+2-1=3$. The formula fits: $U+W=\mathbb{R}^{3}$, and numpy
        gives the rank of the four spanning vectors together as $3$.

  \item Elimination puts pivots in columns $1,2,4$, so the rank is $3$; a basis of the
        column space is $(1,2,0)$, $(2,4,1)$, and $(4,9,1)$. Nullity $4-3=1$ with
        basis $\{(-1,-1,1,0)\}$; \texttt{matrix\_rank} and
        \texttt{null\_space} confirm both numbers.

  \item False. In $\mathbb{R}^{2}$ take $U_{1}=\operatorname{span}\{(1,0)\}$,
        $U_{2}=\operatorname{span}\{(0,1)\}$, $U_{3}=\operatorname{span}\{(1,1)\}$.
        All pairwise intersections are $\{\mathbf{0}\}$, yet
        $(1,1)=(1,0)+(0,1)$: the representation is not unique, and
        $\dim U_{1}+\dim U_{2}+\dim U_{3}=3>2=\dim\mathbb{R}^{2}$.

  \item Take a basis $\mathbf{u}_{1},\dots,\mathbf{u}_{n-1}$ of $U$. That set is
        linearly independent and has fewer than $n$ elements, so it does not yet span
        $V$; there is some $\mathbf{v}\in V$ outside its span. The extension theorem makes
        $\mathbf{u}_{1},\dots,\mathbf{u}_{n-1},\mathbf{v}$ a basis of $V$, so that
        $V=U\oplus\operatorname{span}\{\mathbf{v}\}$.

  \item A polynomial of degree at most $4$ that vanishes at $0$ and $1$ is divisible
        by $x(x-1)$, so $p=x(x-1)(a+bx+cx^{2})$: three free parameters,
        dimension $3$. Basis $\{x(x-1),\,x^{2}(x-1),\,x^{3}(x-1)\}$. The evaluation
        matrix $2\times5$ has rank $2$, so nullity $5-2=3$.

  \item The coordinates satisfy $c_{1}+c_{2}+c_{3}+c_{4}=4$, $c_{2}+c_{3}+c_{4}=2$,
        $c_{3}+c_{4}=-5$, $c_{4}=1$. From the bottom: $c_{4}=1$, $c_{3}=-6$,
        $c_{2}=7$, $c_{1}=2$. So $[q]_{B}=(2,7,-6,1)$, the same as the result of
        \texttt{numpy.linalg.solve}.
\end{enumerate}

\section*{Chapter 7: Basis, Dimension, and Linear Transformations}
\begin{enumerate}[label=\textbf{7.\arabic*}]
  \item Solve $c_{1}(1,1,0)+c_{2}(1,0,1)+c_{3}(0,1,1)=(2,-1,3)$, that is
        $c_{1}+c_{2}=2$, $c_{1}+c_{3}=-1$, $c_{2}+c_{3}=3$. The result is
        $c=(-1,3,0)$, so $[\mathbf{w}]_{B'}=(-1,3,0)$. Back:
        $-1(1,1,0)+3(1,0,1)+0=(2,-1,3)$.

  \item Let $\varphi$ be an isomorphism. Because $\varphi$ is bijective,
        $\mathbf{b}_{i}=\varphi^{-1}(\mathbf{e}_{i})$ is defined; because
        $\varphi^{-1}$ is linear and bijective, $\{\mathbf{b}_{i}\}$ is a basis and
        $\varphi=\varphi_{B}$. It is not unique because each basis $B$ gives a
        different isomorphism.

  \item $A$ has rank $2$ (pivots in columns $2,3$), nullity $1$. Kernel:
        $x_{2}=0$, $2x_{3}=0$, so $\ker A=\operatorname{span}\{(1,0,0)\}$.
        Range $\operatorname{span}\{(1,0,0),(0,2,0)\}$. This is the derivative matrix
        on $P_{2}$ in the basis $(1,x,x^{2})$: $D(1)=0$, $D(x)=1$,
        $D(x^{2})=2x$.

  \item $\operatorname{im}(S\circ T)\subseteq\operatorname{im}S$ gives
        rank$(S\circ T)\le$ rank $S$. Besides that,
        $\ker T\subseteq\ker(S\circ T)$, so nullity$(S\circ T)\ge$ nullity $T$;
        rank-nullity applied to both maps gives rank$(S\circ T)\le$ rank $T$.

  \item False. $\begin{pmatrix}1&0\\0&0\end{pmatrix}$ and
        $\begin{pmatrix}0&1\\0&0\end{pmatrix}$ both have rank $1$, but
        the second is nilpotent (trace $0$, square zero) while the first
        has trace $1$. Rank measures only the dimension of the range.

  \item $T$ sends $\mathbf{b}_{1},\mathbf{b}_{2},\mathbf{b}_{3}$ to $(1,0)$,
        $(3,1)$, $(3,2)$. Coordinates relative to $C$ are $(x,y)\mapsto(y,x-y)$,
        so the columns are $(0,1),(1,2),(2,1)$ and
        $[T]_{C\leftarrow B}=\begin{pmatrix}0&1&2\\1&2&1\end{pmatrix}$.
        For $\mathbf{v}=(1,2,3)$: $[\mathbf{v}]_{B}=(-1,-1,3)$ gives
        $(5,0)$, the same as $[T(\mathbf{v})]_{C}$.

  \item Take a basis $\mathbf{u}_{1},\dots,\mathbf{u}_{m}$ of $U$ with
        $m=\dim V$. It is linearly independent and has $\dim V$ elements, so it spans
        $V$; hence every $\mathbf{v}\in V$ lies in $U$, that is $U=V$.

  \item $B=P^{-1}AP$, so $B\mathbf{x}=\mathbf{0}\iff AP\mathbf{x}=\mathbf{0}
        \iff P\mathbf{x}\in\ker A$. The mapping $\mathbf{x}\mapsto P\mathbf{x}$ is
        an isomorphism from $\ker B$ onto $\ker A$, so the nullities agree. That
        equality is exactly the statement that the dimension of the kernel does not
        depend on the basis.

  \item The third row is the sum of the first two, so rank $2$ and nullity
        $4-2=2$. A nonzero kernel vector $(-2,-1,1,0)$: $F(-2,-1,1,0)=\mathbf{0}$.
        That direction shifts no feature at all, because
        $F(\mathbf{x}+t\mathbf{k})=F(\mathbf{x})$ for all $t$.
\end{enumerate}

\section*{Chapter 8: Linear Systems, Determinants, and Inverses}
\begin{enumerate}[label=\textbf{8.\arabic*}]
  \item The pivots that appear are $1$, $1$, $-1$; no row needs swapping, so
        $\det A=1\cdot1\cdot(-1)=-1$. Backward: $z=3$, $y=-1$, $x=2$. Cramer
        gives $\det A_{x}=-2$, $\det A_{y}=1$, $\det A_{z}=-3$; divided by $-1$
        the result is $(2,-1,3)$.

  \item Cofactor expansion gives $\det M=0$. The refuting vector is
        $\mathbf{v}=(1,-1,-1,1)$: every row gives $0$, for instance
        $0-1-2+3=0$. So $\det M=0$ without any expansion. numpy gives
        $5.3\times10^{-15}$, not zero, because elimination subtracts nearly equal
        numbers and the rounding error is counted too.

  \item Leibniz formula: $\det A=\sum_{\sigma}\operatorname{sgn}(\sigma)
        \prod_{i}a_{i\sigma(i)}$. If columns $j$ and $k$ agree, pair every
        $\sigma$ with $\sigma\circ(j\,k)$; the products of the terms are the same but
        the signs are opposite, so the sum of all terms is zero.

  \item Multiply $A\operatorname{adj}A=(\det A)I$ by the determinant:
        $\det A\cdot\det(\operatorname{adj}A)=(\det A)^{n}$, because
        $\det(cI)=c^{n}$. If $\det A\ne0$, divide by $\det A$. If
        $\det A=0$, both sides vanish for $n\ge2$, so the equality still holds.

  \item False as a general statement. Sensitivity is measured by
        $\kappa(A)=\lVert A\rVert\,\lVert A^{-1}\rVert$, not by $\det A$:
        $A=\varepsilon I$ has $\det A=\varepsilon^{n}$ very small yet
        $\kappa=1$. A small determinant is dangerous only when it comes with a
        large $\lVert A^{-1}\rVert$.

  \item $\lVert A\rVert_{1}=2+\varepsilon$ and
        $A^{-1}=\frac1\varepsilon\begin{pmatrix}1+\varepsilon&-1\\-1&1\end{pmatrix}$
        give $\lVert A^{-1}\rVert_{1}=(2+\varepsilon)/\varepsilon$. So
        $\kappa_{1}(A)=(2+\varepsilon)^{2}/\varepsilon$. The condition $\kappa>10^{6}$
        is $\varepsilon<4.000016\times10^{-6}$; at
        $\varepsilon=4\times10^{-6}$ numpy gives $1.000004\times10^{6}$.

  \item Adding a multiple of a row means multiplying on the left by
        $E=I+cE_{ij}$, and $\det E=1$, so $\det B=\det A$. Because $B=EA$,
        we have $B^{-1}=A^{-1}E^{-1}$ with $E^{-1}=I-cE_{ij}$.

  \item $A+AD=A(I+D)$, and $I+D$ is diagonal, so
        $\det(A+AD)=\det A\prod_{i}(1+d_{i})$. It is invertible exactly when
        $\det A\ne0$ and no $d_{i}=-1$. Example $\det A=5$ with
        $D=\operatorname{diag}(2,-\tfrac12,3)$: the left side is $30$, the right side
        $5\cdot3\cdot\tfrac12\cdot4=30$.

  \item With $A=\begin{pmatrix}1&2&3\\0&1&4\\5&6&0\end{pmatrix}$,
        $\det A=1$ and
        \[
          \operatorname{adj}A=A^{-1}=
          \begin{pmatrix}-24&18&5\\20&-15&-4\\-5&4&1\end{pmatrix},
        \]
        the same as \texttt{numpy.linalg.inv} up to machine error. For
        $\det A=7$:
        $\begin{pmatrix}2&1&0\\1&1&0\\0&0&7\end{pmatrix}^{-1}$ contains
        $1/7$ in the third entry, because division by $7$ cannot be avoided.
\end{enumerate}

\section*{Chapter 9: Eigenvalues and Diagonalization}
\begin{enumerate}[label=\textbf{9.\arabic*}]
  \item $\operatorname{tr}A_{1}=5$ and $\det A_{1}=4$, so $\lambda=4$ and
        $\lambda=1$. Eigenvectors: $(1,1)$ for $\lambda=4$, $(1,-2)$ for
        $\lambda=1$. Then $P=\begin{pmatrix}1&1\\1&-2\end{pmatrix}$ and
        $D=\operatorname{diag}(4,1)$; $P^{-1}A_{1}P=D$ is verified.

  \item $\lambda=5,2$ with eigenvectors $(2,1)$ and $(1,-1)$. For
        $k\ge1$,
        \[
          G^{k}=\frac13
          \begin{pmatrix}2\cdot5^{k}+2^{k}&2\cdot5^{k}-2\cdot2^{k}\\
          5^{k}-2^{k}&5^{k}+2\cdot2^{k}\end{pmatrix},
        \]
        so that
        $G^{10}=\begin{pmatrix}6510758&6509734\\3254867&3255891\end{pmatrix}$,
        the same as \texttt{matrix\_power}.

  \item $A\mathbf{v}=\lambda\mathbf{v}$ gives
        $A^{2}\mathbf{v}=\lambda^{2}\mathbf{v}$; induction gives
        $A^{k}\mathbf{v}=\lambda^{k}\mathbf{v}$. The converse fails:
        $A=\begin{pmatrix}0&-1\\1&0\end{pmatrix}$ has $A^{2}=-I$ with eigenvalue
        $-1$, whereas $A$ itself has no real eigenvalue.

  \item $0$ is an eigenvalue $\iff$ there is $\mathbf{v}\ne\mathbf{0}$ with
        $A\mathbf{v}=\mathbf{0}$ $\iff$ $\ker A\ne\{\mathbf{0}\}$ $\iff$ $A$ is not
        invertible. If $A$ is invertible and $A\mathbf{v}=\lambda\mathbf{v}$, then
        $\lambda\ne0$ and $A^{-1}\mathbf{v}=\lambda^{-1}\mathbf{v}$.

  \item False. $\begin{pmatrix}1&0\\0&1\end{pmatrix}$ and
        $\begin{pmatrix}1&1\\0&1\end{pmatrix}$ both have characteristic
        polynomial $(\lambda-1)^{2}$, but the second is defective:
        $\operatorname{rank}(J-I)=1$, so nullity $1<2$.

  \item $\lambda I-A$ and its transpose have the same determinant, because
        $\det X=\det X^{\top}$; hence $A$ and
        $A^{\top}$ have the same characteristic polynomial and the same eigenvalues.
        Their eigenvectors do not: $A$ and $A^{\top}$ have different eigenspaces as
        soon as $A$ is not symmetric.

  \item $K-2I=\begin{pmatrix}0&1&0\\0&0&0\\0&0&-1\end{pmatrix}$ has rank $2$,
        so nullity $=1$. The algebraic multiplicity of $\lambda=2$ is $2$,
        the geometric one is $1$; because $1<2$, $K$ cannot be diagonalized.

  \item $\lambda=1$ and $0.5$. The fixed point of $x=4y$ with $x+y=1$ is
        $(0.8;0.2)$. The error shrinks like $0.5^{k}$; because
        $0.5^{19}=1.91\times10^{-6}>\allowbreak10^{-6}$ and
        $0.5^{20}=9.54\times10^{-7}$, it takes $20$ steps.

  \item The printed matrix gives
        $\operatorname{tr}(B^{10})=(-19)+21+1024=1026$, and
        $1^{10}+1^{10}+2^{10}=1026$: they agree. For a diagonalizable $A$,
        $A^{k}=PD^{k}P^{-1}$ yields
        $\operatorname{tr}(A^{k})=\operatorname{tr}(D^{k})=\sum_{i}\lambda_{i}^{k}$.
        For defective $B$ that step is unavailable, because $B$ has no $P$ and $D$.
        Instead use the Schur triangular form $B=QTQ^{-1}$ with $T$ upper triangular
        and diagonal $\lambda_{i}=1,1,2$; then $B^{k}=QT^{k}Q^{-1}$, and $T^{k}$ is
        still upper triangular with diagonal $\lambda_{i}^{k}$, so
        $\operatorname{tr}(B^{k})=\operatorname{tr}(T^{k})
        =\sum_{i}\lambda_{i}^{k}=1026$. The trace is invariant under change of basis
        ($\operatorname{tr}XY=\operatorname{tr}YX$), which is why the result does not
        depend on $Q$. In short: diagonalization carries the proof only in the
        diagonalizable case; in the defective case the proof switches to the
        triangular form, the formula stays the same, and the factor $k$ that appears
        in the entries of $B^{k}$ never touches the diagonal.
\end{enumerate}

\section*{Chapter 10: Orthogonality, SVD, and Least Squares}
\begin{enumerate}[label=\textbf{10.\arabic*}]
  \item $(1,0,0)$ and $(0,1,1)$ are already orthogonal, so
        $\operatorname{proj}_{W}\mathbf{b}=1\cdot(1,0,0)+\frac{4}{2}(0,1,1)
        =(1,2,2)=\mathbf{b}$. Because $\mathbf{b}\in W$, the distance is $0$;
        \texttt{qr} gives $QQ^{\top}\mathbf{b}=(1,2,2)$ as well.

  \item $\mathbf{b}=\mathbf{p}+\mathbf{r}$ with $\mathbf{r}\perp W$, so
        $\mathbf{r}\perp\mathbf{p}$ and Pythagoras gives
        $\lVert\mathbf{b}\rVert^{2}=\lVert\mathbf{p}\rVert^{2}
        +\lVert\mathbf{r}\rVert^{2}$. If $\mathbf{p}$ is not an orthogonal
        projection, $\mathbf{r}$ is not perpendicular to $\mathbf{p}$; on an oblique
        basis in $\mathbb{R}^{2}$ a cross term $2\langle\mathbf{p},\mathbf{r}
        \rangle\ne0$ appears.

  \item Orthonormal columns: $(A^{\top}A)_{ij}=\mathbf{a}_{i}^{\top}\mathbf{a}_{j}
        =\delta_{ij}$, so $A^{\top}A=I$. The normal equations
        $A^{\top}A\hat{\mathbf{x}}=A^{\top}\mathbf{b}$ immediately give
        $\hat{\mathbf{x}}=A^{\top}\mathbf{b}$. All singular values of $A$ are
        $1$, so $\kappa(A)=1$.

  \item False. $A=\begin{pmatrix}1&0\\0&1\\1&1\end{pmatrix}$ gives
        $A^{\top}A=\begin{pmatrix}2&1\\1&2\end{pmatrix}$, which is invertible
        ($\det=3$), yet its rows are dependent: the third row is the sum of the first
        two. What invertibility of $A^{\top}A$ guarantees is only independence of the
        columns.

  \item Singular values $2$ and $0$, so rank $1$;
        $U=V=\frac{1}{\sqrt2}\begin{pmatrix}1&1\\1&-1\end{pmatrix}$ (up to
        sign). $A^{+}=\frac14\begin{pmatrix}1&1\\1&1\end{pmatrix}$, so
        $A^{+}\mathbf{b}=(1,1)$ and the residual is $(-1,1)$ with norm
        $\sqrt2=1.4142$; \texttt{pinv} agrees.

  \item 
Write $\mathbf{b}=A\hat{\mathbf{x}}+\mathbf{r}$ with
        $A\hat{\mathbf{x}}\in\operatorname{col}A$. For
        $\mathbf{h}\in\operatorname{col}A$ the function
        $t\mapsto\lVert\mathbf{b}-A(\hat{\mathbf{x}}+t\mathbf{h})\rVert^{2}$ is
        minimal at $t=0$, hence $\langle\mathbf{r},\mathbf{h}\rangle=0$; this is the
        ``necessary'' direction. Conversely, if
        $\langle\mathbf{r},\mathbf{h}\rangle=0$ for every
        $\mathbf{h}\in\operatorname{col}A$, then
        $\lVert\mathbf{b}-A\mathbf{x}\rVert^{2}=\lVert\mathbf{r}\rVert^{2}+\lVert
        A(\hat{\mathbf{x}}-\mathbf{x})\rVert^{2}\ge\lVert\mathbf{r}\rVert^{2}$, so
        $\hat{\mathbf{x}}$ minimises. Column independence is what makes
        $\hat{\mathbf{x}}$ unique.

  \item $\mathbf{u}_{1}=(1,1,1)/\sqrt3$; the remainder $(1,1,0)-\frac23(1,1,1)
        =\frac13(1,1,-2)$, so $\mathbf{u}_{2}=(1,1,-2)/\sqrt6$; then
        $\mathbf{u}_{3}=(1,-1,0)/\sqrt2$. $\lVert Q^{\top}Q-I\rVert$:
        $6.7\times10^{-16}$ (by hand), $3.1\times10^{-16}$ (\texttt{qr}).
        The signs may differ because each $\mathbf{u}_{i}$ is determined only up to
        $\pm$.

  \item $AA^{\top}\mathbf{x}=\mathbf{0}$ gives
        $\lVert A^{\top}\mathbf{x}\rVert^{2}=\mathbf{x}^{\top}AA^{\top}
        \mathbf{x}=0$, so $A^{\top}\mathbf{x}=\mathbf{0}$; independence of the rows
        makes $\mathbf{x}=\mathbf{0}$. Hence $AA^{\top}$ is square and invertible.
        For $\mathbf{x}=A^{\top}(AA^{\top})^{-1}\mathbf{b}$ we have
        $A\mathbf{x}=\mathbf{b}$ and
        $\mathbf{x}\in\operatorname{col}A^{\top}=(\ker A)^{\perp}$.

  \item With the last reading $314.3$: $\sum y=1540.5$,
        $\sum Ty=157161.7$, determinant $50$. Cramer:
        $a=(52030\cdot1540.5-510\cdot157161.7)/50=-5.04$ and
        $b=(5\cdot157161.7-510\cdot1540.5)/50=3.07$;
        \texttt{lstsq} agrees. The intercept shifts by $-10.10$, the slope by $+0.10$;
        the endpoint has leverage $0.6$, the largest of the five.

  \item $B$ has rank $1$, so $B^{+}$ gives
        $\mathbf{x}^{+}=\frac{1}{70}(1,2)\approx(0.01429;\,0.02857)$ and
        $\mathbf{k}=(-2,1)$. All solutions are $\mathbf{x}^{+}+t\mathbf{k}$;
        because $\mathbf{x}^{+}\perp\mathbf{k}$,
        $\lVert\mathbf{x}^{+}+t\mathbf{k}\rVert^{2}=\lVert\mathbf{x}^{+}
        \rVert^{2}+t^{2}\lVert\mathbf{k}\rVert^{2}$ is minimal at $t=0$.
\end{enumerate}

\section*{Chapter 11: Continuity and Uniform Continuity}
\begin{enumerate}[label=\textbf{11.\arabic*}]
  \item $\lvert(2x-1)-5\rvert=2\lvert x-3\rvert$, so $\delta=\varepsilon/2$ is
        valid and largest: at $x=3+\delta$ the value is exactly $\varepsilon$.
        The choice $\delta=\varepsilon/3$ is also valid, but it uses only a third of the
        room available.

  \item $\lvert x^{2}-c^{2}\rvert=\lvert x-c\rvert\,\lvert x+c\rvert$; with
        $\lvert x-c\rvert<1$ we get $\lvert x+c\rvert\le2\lvert c\rvert+1$,
        so $\delta=\min(1,\varepsilon/(2\lvert c\rvert+1))$. At $c=2$:
        $\delta=0.002$; at $c=100$: $\delta=4.975\times10^{-5}$.

  \item Partly true. $1/x$ is indeed continuous at every point of $(0,1]$. But the
        uniform conclusion is wrong: $(0,1]$ is not compact and $1/x$ is unbounded
        near $0$. The refuting pair $x_{n}=1/n$, $y_{n}=1/(n+1)$:
        $\lvert x_{n}-y_{n}\rvert\to0$ but the difference of the values is $1$.

  \item Given $\varepsilon$, take $\delta_{f}$ and $\delta_{g}$; with
        $\delta=\min(\delta_{f},\delta_{g})$,
        $\lvert(f+g)(x)-(f+g)(y)\rvert\le\varepsilon/2+\varepsilon/2$.
        For $fg$: $f=g=x$ on $\mathbb{R}$ with $x_{n}=n$,
        $y_{n}=n+1/n$ gives difference $2+1/n^{2}\ge2$.

  \item $\lvert\sqrt x-\sqrt y\rvert^{2}\le\lvert x-y\rvert$ (rationalize the
        denominator), so $\delta=\varepsilon^{2}$ works for all
        $x,y\ge0$. Only $x,y\ge0$ is used, so that the roots are defined;
        closedness is not needed because the proof never touches a limit point outside
        the domain.

  \item 
$(1+\delta)^{3}-1=\varepsilon$ gives
        $\delta=(1+\varepsilon)^{1/3}-1=0{,}0033223$. The choice
        $\delta=\varepsilon/3=0{,}0033333$ is actually too large:
        $(1+\varepsilon/3)^{3}-1-\varepsilon=\varepsilon^{2}/3+\varepsilon^{3}/27>0$
        for every $\varepsilon>0$, so it violates the bound and is not valid.

  \item With $\delta=\varepsilon/2$: the largest difference
        $2\delta-\delta^{2}=0.009975<\varepsilon$, so it is safe but not
        optimal. With $\delta=1-\sqrt{1-\varepsilon}=0.0050126$: the largest
        difference $1-(1-\delta)^{2}=\varepsilon$ exactly, attained at
        $x=1-\delta$.

  \item $\lvert f\rvert$ is continuous on the compact $[a,b]$, so it attains a minimum $m$;
        because $f\ne0$ everywhere, $m>0$ and $\lvert f\rvert\ge m$. For
        $f(x)=x$ on $(0,1)$ the conclusion fails: the infimum is $0$ and it is not
        attained.

  \item At $c\ne0$ take $q_{n}\to c$ rational and $r_{n}\to c$ irrational; the
        function values tend to $c$ and $-c$, which differ, so $f$ is not continuous. At $0$:
        $\lvert f(x)\rvert=\lvert x\rvert\to0=f(0)$. The function $\lvert f\rvert$
        is $x\mapsto\lvert x\rvert$, continuous everywhere; there is no
        contradiction because $\lvert f\rvert$ throws away the sign information.
\end{enumerate}

\section*{Chapter 12: Derivatives: One and Several Variables}
\begin{enumerate}[label=\textbf{12.\arabic*}]
  \item $\bigl((2+h)^{3}-8\bigr)/h=12+6h+h^{2}\to12$. For $\sqrt x$:
        $(\sqrt{2+h}-\sqrt2)/h=1/(\sqrt{2+h}+\sqrt2)\to1/(2\sqrt2)
        =0.3535534$. In both cases $h$ drops out of the denominator.

  \item $f'(c)>0$ makes $(f(c+h)-f(c))/h>f'(c)/2>0$ for $h$ small enough,
        so $f(c+h)>f(c)$. That is only a one-sided statement. The example
        $f(x)=x+2x^{2}\sin(1/x^{2})$ has $f'(0)=1>0$ but is not increasing on
        any neighbourhood of $0$: its derivative changes sign arbitrarily close
        to $0$.

  \item Differentiate $1/g$ first: $(1/g)'=-g'/g^{2}$; here the continuity of $g$
        is used, so that $g(x+h)\to g(x)\ne0$. Then
        $(f/g)'=f'\cdot(1/g)+f\cdot(1/g)'$, which simplifies to
        $(f'g-fg')/g^{2}$.

  \item False. $f(x,y)=xy/(x^{2}+y^{2})$ for $(x,y)\ne(0,0)$ and
        $f(0,0)=0$: both partial derivatives at the origin exist and equal $0$,
        but along $y=x$ the value is $\tfrac12$, so $f$ is not continuous
        there.

  \item $\nabla f=(3x^{2}y-y^{3},\,x^{3}-3xy^{2})$, at $(2,1)$ equal to
        $(11,2)$. The steepest ascent direction is $(11,2)/\sqrt{125}$, the maximum rate
        $\sqrt{125}=11.1803$. The central difference with $h=10^{-4}$ gives relative error
        $9.1\times10^{-10}$ and $1.0\times10^{-8}$, consistent with order $h^{2}$ in the
        second component.

  \item By the mean value theorem,
        $f(b)-f(a)=f'(\xi)(b-a)=0$ for every pair $a<b$ in the interval, so
        $f$ is constant. On a disconnected set that step fails because $\xi$
        need not lie in the set; for example $f=0$ on $(0,1)$ and $f=1$
        on $(2,3)$.

  \item At $(0,1)$: $J_{H}=\begin{pmatrix}1&1\\1&2\end{pmatrix}$; at
        $H(0,1)=(1,1)$: $J_{G}=\begin{pmatrix}2&1\\1&-1\end{pmatrix}$. The matrix
        product is $\begin{pmatrix}3&4\\0&-1\end{pmatrix}$, and the direct derivative
        gives the same matrix.

  \item $f_{x}=3x^{2}-3y$, $f_{y}=3y^{2}-3x$; both zero gives $y=x^{2}$ and
        $x=y^{2}$, so $x(x^{3}-1)=0$: the points $(0,0)$ and $(1,1)$. The Hessian
        $\begin{pmatrix}6x&-3\\-3&6y\end{pmatrix}$: at $(0,0)$ its determinant is
        $-9<0$ (saddle), at $(1,1)$ its determinant is $27>0$ with $f_{xx}=6>0$
        (local minimum, $f=-1$). If the determinant is zero, that test says nothing.

  \item $p_{2}(x)=1+\tfrac12h-\tfrac18h^{2}$ at $h=0.04$ gives $1.0198$,
        whereas $\sqrt{1.04}=1.0198039$; the error is
        $3.9\times10^{-6}$, the Lagrange remainder
        $\frac{3}{8}\xi^{-5/2}h^{3}/6\le4.0\times10^{-6}$. The bound is tight, not loose, because $\xi$ is close
        to $1$: the Lagrange remainder really does nearly equal the actual error.
\end{enumerate}

\section*{Chapter 13: Riemann Integrals and the Fundamental Theorem of Calculus}
\begin{enumerate}[label=\textbf{13.\arabic*}]
  \item $U_{n}=\frac1n\sum_{k=1}^{n}\frac kn=\frac{n+1}{2n}$ and
        $L_{n}=\frac{n-1}{2n}$, so $U_{n}-L_{n}=1/n\to0$: $f$ is integrable
        and $\int_{0}^{1}x\,dx=\tfrac12$.

  \item $\bigl\lvert\lvert f(x)\rvert-\lvert f(y)\rvert\bigr\rvert
        \le\lvert f(x)-f(y)\rvert$ (reverse triangle inequality), so
        integrability carries over. The bound comes from $-\lvert f\rvert\le f\le\lvert f\rvert$
        and monotonicity of the integral. The converse is false: $f=1$ on the rationals and
        $-1$ on the irrationals has $\lvert f\rvert=1$, which is integrable.

  \item False. The correct part: the upper and lower sums of a bounded function are always
        finite. What the proof does not support is the equality
        $\sup_{\mathcal{P}}L=\inf_{\mathcal{P}}U$. The Dirichlet function is bounded
        with $U-L=1$ for every partition.

  \item On a uniform partition each subinterval contributes oscillation
        $f(x_{k})-f(x_{k-1})$, so that
        $U-L\le\frac{b-a}{n}\sum_{k}\bigl(f(x_{k})-f(x_{k-1})\bigr)
        =\frac{(b-a)(f(b)-f(a))}{n}$. For a non-monotone function that telescoping fails
        because the signs of the oscillations are not uniform.

  \item $\int_{0}^{1}(3x^{2}-2x+1)\,dx=[x^{3}-x^{2}+x]_{0}^{1}=1$. Simpson
        with $n=100$ gives $1.0000000000$ and \texttt{quad} $1.0000000000$
        (error $1.1\times10^{-14}$); the three numbers differ by zero to machine
        precision.

  \item $\int_{1}^{10}x^{-2}dx=1-\tfrac1{10}=0.9$ and
        $\int_{1}^{\infty}x^{-2}dx=1$. Their difference is $0.1=\tfrac1{10}$ exactly,
        namely the tail $\int_{10}^{\infty}x^{-2}dx=1/10$; there is no need to recompute it,
        because the missing part is exactly that tail.

  \item If $f(x_{0})>0$, continuity gives $f>\tfrac12f(x_{0})$ on an interval of
        width $2\delta$ around $x_{0}$ (of width $\delta$ when $x_{0}$ is an
        endpoint), so the integral is at least
        $\tfrac12\delta f(x_{0})>0$: a contradiction. For $f$ that is merely integrable
        the statement fails: $f=1$ at a single point and $0$ elsewhere has integral
        $0$.

  \item $u=\sin x$: $\int_{0}^{1}u\,du=\tfrac12$. Through the identity:
        $\tfrac12\int_{0}^{\pi/2}\sin2x\,dx
        =\tfrac12\bigl[-\tfrac12\cos2x\bigr]_{0}^{\pi/2}=\tfrac12$. \texttt{quad}
        gives $0.5$ as well.

  \item On every subinterval there is a rational (value $1$) and an irrational (value
        $0$), so $U(\mathcal{P},D)=1$ and $L(\mathcal{P},D)=0$; the difference is $1$
        for all $\mathcal{P}$. The Riemann criterion demands that this difference can be
        made as small as we please, so $D$ is not integrable.

  \item The integral bounds give $\frac1{101}<\sum_{n=101}^{\infty}n^{-2}
        <\frac1{100}$, that is $9.9010\times10^{-3}$ up to
        $1.0000\times10^{-2}$. Their midpoint is $9.9505\times10^{-3}$,
        whereas the true value is $9.950167\times10^{-3}$; the relative error is
        $3.3\times10^{-5}$.
\end{enumerate}

\section*{Chapter 14: Series of Functions and the Exchange of Limits}
\begin{enumerate}[label=\textbf{14.\arabic*}]
  \item $\sup_{[0,1]}\lvert x^{n}\rvert=1$ for all $n$, attained at $x=1$.
        At $x=0.9$: $0.3487$ for $n=10$, $2.656\times10^{-5}$ for
        $n=100$, $1.748\times10^{-46}$ for $n=1000$. The first column is flat
        because the peak is always at the endpoint; the second column shrinks because $0.9<1$.

  \item For $n\ge N$, $\lVert f_{n}-f\rVert_{\infty}<\varepsilon$, so
        $\lVert f\rVert_{\infty}\le\lVert f_{N}\rVert_{\infty}+1$ and likewise
        for $\lVert f_{n}\rVert_{\infty}$. Pointwise convergence is not enough:
        $f_{n}=n\chi_{(0,1/n)}$ converges to $0$ at every point but
        $\lVert f_{n}\rVert_{\infty}=n$.

  \item The peak is at $x=1/\sqrt n$ with
        $\lVert f_{n}\rVert_{\infty}=1/(2\sqrt n)\to0$, and that bound holds
        for all $x$ at once, so the convergence is uniform. The condition
        $1/(2\sqrt n)<0.01$ requires $\sqrt n>50$, so $n=2501$.

  \item False. Pointwise reasoning gives one $N$ for each $x$, not
        one $N$ for all $x$. The counterexample is $x_{n}=2^{-1/n}$:
        $x_{n}^{n}=\tfrac12$ for every $n$, so
        $\sup_{[0,1)}\lvert x^{n}\rvert\ge\tfrac12$ and it does not tend to $0$.

  \item The error $\lvert e-S_{N}(1)\rvert$ for $N=5,10,15$:
        $1.615\times10^{-3}$, $2.731\times10^{-8}$,
        $5.077\times10^{-14}$; the bound $e/(N+1)!$ gives
        $3.775\times10^{-3}$, $6.810\times10^{-8}$,
        $1.299\times10^{-13}$, always larger. At $N=20$ the bound is already
        $5.32\times10^{-20}$, far below machine precision, so the measured error
        stops at $4.44\times10^{-16}$. The bound first drops below $10^{-16}$
        at $N=18$, and there the true error is unmeasurable as well.

  \item The sum is $(1-x)/(1-x)=1$ for $x\in[0,1)$, whereas at $x=1$ the
        sum is $0$. The limit is not continuous even though every partial sum is
        continuous; the uniform limit of continuous functions must be continuous, so
        the convergence is not uniform.

  \item Radii: $R=1$ for $\sum x^{n}/n^{2}$, $R=e$ for
        $\sum(n!/n^{n})x^{n}$, $R=\tfrac12$ for $\sum2^{n}x^{n}/n$. The ratio of the
        $1001$st and $1000$th terms: $0.998003$; $0.368063$ (so
        $R\approx2.7169$); $1.998002$. The slowest is $\sum x^{n}/n^{2}$,
        because its ratio tends to $1$ like $(1-1/n)^{2}$.

  \item $\lvert f_{n}g-fg\rvert=\lvert f_{n}-f\rvert\,\lvert g\rvert
        \le\lVert f_{n}-f\rVert_{\infty}\lVert g\rVert_{\infty}\to0$. Without
        boundedness of $g$: $f_{n}=x+1/n\to x$ uniformly on $(0,1)$, but
        $(x+1/n)/x=1+1/(nx)$ is unbounded, so $f_{n}g$ does not converge
        uniformly.

  \item The peak is at $x=1/\sqrt{2n+1}$. For $n=10$: height $1.3397$ at
        $x=0.2182$, $\int_{0}^{1}f_{10}=\tfrac{5}{11}=0.4545$. For
        $n=100$: height $4.2835$ at $x=0.07053$, the integral
        $\tfrac{50}{101}=0.4950$. The height follows $\sqrt{n/(2e)}$; the area
        tends to $\tfrac12$ because the peak gets narrower.

  \item Differentiating $\sum_{n\ge0}x^{n}=1/(1-x)$ gives
        $\sum_{n\ge1}nx^{n-1}=1/(1-x)^{2}$; at $x=\tfrac12$ we get
        $\sum_{n\ge1}n/2^{n}=2$. The tail remainder is $(N+2)/2^{N}$, so
        $N=39$ gives $7.46\times10^{-11}<10^{-10}$.
\end{enumerate}

\section*{Chapter 15: Multiple Integrals, Change of Variables, and Vector Field Theorems}
\begin{enumerate}[label=\textbf{15.\arabic*}]
  \item Cartesian: $\int_{-1}^{1}\int_{-\sqrt{1-x^{2}}}^{\sqrt{1-x^{2}}}
        (x^{2}+y^{2})\,dy\,dx=\pi/2$. Polar:
        $\int_{0}^{2\pi}\int_{0}^{1}r^{3}\,dr\,d\theta
        =2\pi\cdot\tfrac14=\pi/2=1.5708$. \texttt{dblquad} gives
        $1.5707963268$; in this test environment the Cartesian route requires $8379$
        evaluations, the polar route $441$.

  \item For $T(\mathbf{u})=A\mathbf{u}$ linear, $DT=A$ is constant, so
        its Jacobian determinant is $\det A$ and $\lvert\det A\rvert$ is the area
        scale factor. If $\det A\ne0$ the map is one-to-one on $D$, so
        $\iint_{T(D)}f=\lvert\det A\rvert\iint_{D}f\circ T$.

  \item Clockwise, each side is traversed in reverse, so the four terms of the example
        change sign: $-\tfrac13-2-\tfrac23+1=-2$, whereas the right-hand side stays
        $2$. The bottom and right sides that used to contribute $+\tfrac13$ and $+2$
        become negative; the left side that used to be $-1$ becomes $+1$.

  \item $x=au$, $y=bv$, $z=cw$ gives Jacobian $abc$ and the unit ball:
        $V=abc\cdot\frac{4\pi}{3}=\frac{4\pi abc}{3}$. At $a=b=c=1$
        the result is $4.1888$, the familiar formula for the volume of a ball.

  \item False. $\operatorname{curl}F=0$ does hold at every point where $F$
        is defined, but the domain has a hole so the curl test does not apply.
        The supporting number: $\oint_{\lvert\mathbf{r}\rvert=1}F\cdot d\mathbf{r}
        =2\pi\ne0$, so $F$ is not conservative.

  \item If the integral does not depend on the path, two paths $C_{1},C_{2}$ from
        $A$ to $B$ give the same integral, so the circulation around $C_{1}-C_{2}$ is zero.
        Conversely, if every closed curve has zero circulation, $C_{1}-C_{2}$ is
        closed and hence the two integrals agree.

  \item $\operatorname{curl}F=(0,0,2)$; the flux through $z=1-x^{2}-y^{2}$ with
        upward normal $=2\cdot\pi=2\pi$. The line integral along the boundary
        $x^{2}+y^{2}=1$, $z=0$:
        $\int_{0}^{2\pi}(-\sin t,\cos t,0)\cdot(-\sin t,\cos t,0)\,dt=2\pi$.
        Both are $6.2832$.

  \item $\operatorname{div}F=0$ away from the origin, so the multiple integral is $0$,
        whereas the flux through the unit sphere is $4\pi=12.5664$. Gauss's theorem
        requires a $C^{1}$ field on a region that contains the whole interior of the
        surface; that condition is violated at the enclosed origin.

  \item $\det J=r^{2}\sin\varphi$, zero at $\varphi=0$ and $\varphi=\pi$. At both
        poles the map stops being one-to-one: all $\theta$ give the same point.
        On $0<\varphi<\pi$ the map is one-to-one and $\det J>0$, so the change of
        variables formula remains valid even though the determinant is zero on the boundary.
\end{enumerate}
